%% kapa: exam=1 round=34 year=2023 period=1 form=A subject=economics subject_ko=경제학원론 seq=26 number=66
%% source: Q-Net 감정평가사 기출문제 (cst00309) artlSeq=5211789 file=raw/1008/5211789/2023년도 제34회 감정평가사 1차시험 1교시.pdf
%% conversion: kordoc 4.24.51 -> tools/gen1.py (auto); review=proofread-sonnet; flags=none
\begin{question}{66}
본원통화에 관한 설명으로 옳지 않은 것은?
\begin{choices}
  \item 본원통화는 현금통화와 은행의 지급준비금의 합과 같다.
  \item 본원통화는 중앙은행의 화폐발행액과 은행의 중앙은행 지급준비예치금의 합과 같다.
  \item 중앙은행의 대차대조표상의 순대정부대출이 증가하면 본원통화는 증가한다.
  \item 중앙은행의 대차대조표상의 순해외자산이 증가하면 본원통화는 증가한다.
  \item 추가로 발행된 모든 화폐가 은행의 시재금(vault cash)으로 보관된다면 본원통화는 증가하지 않는다.
\end{choices}
\end{question}
\answer{5}
